\subsection{Equicontinuity and Stone-Weierstrass}
why are we defining algebras now lmfao

\begin{defn}
	we say $A\subseteq C_{0}(M)$ \textbf{separates pts} if $\fa p,q\in M$, $\ex f\in A$ s.t.
	$f(p)\neq f(q)$.

	we say $A$ \textbf{vanishes} at $x\in M$ if $f(x)=0$ $\fa f\in A$.
\end{defn} \

\begin{thm}[title=Stone-Weierstrass]
	if $M$ cpt, $A\subseteq C_{0}(M)$ an alg of fn's which separates pts and
	vanishes nowhere, then $\bar{A}=C_{0}(M)$.
\end{thm}

\lecdate{Lec 11 - Feb 12 (Week 6)}

before we prove stone-weierstrass, we need the following lemma.
\begin{lm}
	if $A\subseteq C_{0}(M)$ an alg, then so is $\bar{A}$.
\end{lm}

\newpage
\begin{pf}[source=Primary Source Material]
	we'll just prove $\bar{A}$ closed under products. let $f,g\in\bar{A}$.

	then $\ex f_{n},g_{n}\in A$ s.t.
	$f_{n}\rightrightarrows f$ and $g_{n}\rightrightarrows g$.
	then $f_{n}g_{n}\rightrightarrows fg$, so $fg\in\bar{A}$.
\end{pf} \

\begin{lm}
	if $A\subseteq C_{0}(M)$ an alg which separates pts + vanishes nowhere,
	then $\fa p_{1}\neq p_{2}\in M$ and $c_{1},c_{2}\in\bR$,
	$\ex H\in A$ s.t. $H(p_{1})=c_{1}$ and $H(p_{2})=c_{2}$.
\end{lm} \

\begin{pf}[source=Primary Source Material]
	since $A$ vanishes nowhere, $\ex g_{1},g_{2}$ s.t. $g_{1}(p_{1})\neq0$
	and $g_{2}(p_{2})\neq0$. then $g=g_{1}^{2}+g_{2}^{2}$ doesn't vanish at
	either $p_{i}$.

	let $f\in A$ s.t. $f(p_{1})\neq f(p_{2})$. consider:
	\begin{equation*}
		M =
		\begin{pmatrix}
			g(p_{1}) & f(p_{1})g(p_{1}) \\
			g(p_{2}) & f(p_{2})g(p_{2})
		\end{pmatrix} \qquad
		\det(M) = g(p_{1})g(p_{2})(f(p_{2})-f(p_{1})) \neq 0
	\end{equation*}
	therefore, the eq'n:
	\begin{equation*}
		M
		\begin{pmatrix}
			\xi \\ \eta
		\end{pmatrix}
		\ = \
		\begin{pmatrix}
			c_{1} \\ c_{2}
		\end{pmatrix}
	\end{equation*}
	has a soln $(\xi,\eta)$.
	thus, $H=\xi g+\eta fg\in A$ satisfies $H(p_{1})=c_{1},H(p_{2})=c_{2}$.
\end{pf}

we now prove stone-weierstrass. dont feel like boxing it

first, notice if $f\in\bar{A}$, then $\abs{f}\in\bar{A}$.
to see this, suffices to find $f_{n}\in\bar{A}$ s.t. $f_{n}\rightrightarrows\abs{f}$.
consider:
\begin{equation*}
	I \ = \ [-\norm{f}_{\infty},\norm{f}_{\infty}]
\end{equation*}
by weierstrass approx. thm, $\ex p_{n}(x)$ polys s.t. $p_{n}\rightrightarrows\abs{x}$ on $I$.
consider $f_{n}=p_{n}(f)$. since $\bar{A}$ an alg, $f_{n}\in\bar{A}$ for all $n$.
on the other hand, since $p_{n}\rightrightarrows\abs{x}$ on $I$, then
$f_{n}\rightrightarrows\abs{f}$ on $M$.

next, if $f,g\in\bar{A}$, then:
\begin{equation*}
	\max(f,g) \ = \ \frac{f+g}{2}+\frac{\abs{f-g}}{2} \ \in \ \bar{A} \qquad
	\min(f,g) \ = \ \frac{f+g}{2}-\frac{\abs{f-g}}{2} \ \in \ \bar{A}
\end{equation*}
finally, max/min of finitely many elems of $\bar{A}$ are also in $\bar{A}$.
we now prove the statement.

given $F\in C_{0}(M)$ and $\ep>0$, want to find $G\in\bar{A}$ s.t. $\fa x\in M$:
\begin{equation*}
	F(x)-\ep \ \leq \ G(x) \ \leq \ F(x)+\ep
\end{equation*}
by lemma, $\fa p,q\in M$, $\ex H_{pq}\in\bar{A}$ s.t.:
\begin{equation*}
	H_{pq}(p)=F(p) \qquad H_{pq}(q)=F(q)
\end{equation*}
fix $p\in M$. then $\fa q\in M$, $\ex U_{q}\ni q$ nbhd s.t. $\fa x\in U_{q}$:
\begin{equation*}
	H_{pq}(x) \ > \ F(x)-\ep
\end{equation*}
now, $\set{U_{q}}$ set of nbhds arnd every pt $q$ forms an open cvr of $M$.
since $M$ cpt, then $\ex\set{U_{q_{1}},\dots,U_{q_{n}}}$ cvring $M$. consider:
\begin{equation*}
	G_{p}(x) \ = \ \max_{j}H_{pq_{j}}(x)
\end{equation*}
note $G_{p}\in\bar{A}$, and $G_{p}(x)>F(x)-\ep$ for all $x\in M$.
the first inequality has been satisfied.

now, $\fa p\in M$, $\ex$ nbhd $V_{p}\ni p$ s.t. $\fa x\in V_{p}$:
\begin{equation*}
	G_{p}(x)<F(x)+\ep
\end{equation*}
as before, $V_{p_{1}},\dots,V_{p_{m}}$ cvr $M$. let:
\begin{equation*}
	G \ = \ \min_{j}G_{p_{j}}\in\bar{A}
\end{equation*}
similar to before, we get that $G(x)<F(x)+\ep$ for all $M$.
OTOH, $G(x)>F(x)-\ep$ for all $M$, so we have found our desired function.
\qed

we need to patch a small hole near the start:
in particular, $\bar{A}$ may not contain any constant fn's.
\begin{lm}
	sps $I=[-M,M]$. $\ex$ seq of polys w/ no constant term s.t.
	$p_{n}\rightrightarrows\abs{x}$ in $I$.
\end{lm} \

\begin{pf}[source=Primary Source Material]
	by weierstrass approx. thm, $\ex q_{n}(x)=a_{n}+p_{n}(x)$ s.t.
	$q_{n}\rightrightarrows\abs{x}$. then:
	\begin{equation*}
		q_{n}(0)\sto\abs{0}=0 \ \implies \ a_{n}=q_{n}(0)\sto0 \trm{ as } n\sto\infty
	\end{equation*}
	since $a_{n}\sto0$, then $p_{n}=q_{n}-a_{n}\rightrightarrows\abs{x}$.
\end{pf}

preview of a later topic
\begin{defn}
	we say $f$ is a \textbf{trigonometric polynomial} if:
	\begin{equation*}
		f(x)=a_{0}+\sum_{k=1}^{n}a_{k}\cos(kx)+\sum_{k=1}^{m}b_{k}\sin(kx)
	\end{equation*}
	as finite sums.
\end{defn} \

\begin{thm}
	let $f$ be $2\pi$-periodic cts.
	then $\ex p_{n}$ seq of trig polys s.t. $p_{n}\rightrightarrows f$.
\end{thm}
want to use stone-weierstrass, so want to restrict to $[0,2\pi]$.
however, 0 and $2\pi$ are not separated, and we can't exclude either pt while
maintaining cptness. idea: we identify them together instead.

\newpage
\begin{pf}[source=Primary Source Material]
	Identify the space of $2\pi$-periodic cts $f:\bR\gto\bR$ w cts
	$\tilde{f}:S^{1}\gto\bR$ by setting $\tilde{f}(e^{i\theta})=f(\theta)$ for $\theta\in[0,2\pi]$.

	need to check trig polys satisfy hypotheses as elems of $C_{0}(S^{1})$.
	indeed, they're an alg bc of trig identities, nowhere vanishing b/c they
	include constant fn's, and they separate pts (exercise).

	thus, by stone-weierstrass, we're done.
\end{pf}
